% XeLaTeX + ctex 中文支持
\documentclass[UTF8,a4paper,10pt,oneside]{ctexart}

% ---------- 页面与基础 ----------
\usepackage{geometry} % 页面尺寸与页边距设置
\geometry{margin=2.5cm}
\usepackage{indentfirst} % 使章节后的首段缩进
\setlength{\parindent}{2em}
\usepackage{placeins} % 提供 \FloatBarrier 控制浮动体位置

% ---------- 数学与公式 ----------
\usepackage{amsmath,amssymb} % 高级数学公式环境与数学符号

% ---------- 图形与颜色 ----------
\usepackage[table]{xcolor} % 颜色支持（含表格着色）
\usepackage{graphicx} % 插入图片命令 \includegraphics
\usepackage{tikz} % 矢量绘图，配合 pgfplots 使用
\usetikzlibrary{positioning, calc, shapes.geometric, arrows.meta, backgrounds, external, fit} % TikZ 常用库

% ---------- 表格 ----------
\usepackage{booktabs} % 专业表格线命令（\toprule/\midrule/\bottomrule）
\usepackage{tabularx} % 可伸缩列宽的表格环境
\usepackage{array} % 表格列格式增强
\usepackage{multirow} % 合并多行单元格
\usepackage{longtable} % 跨页长表格
\usepackage{colortbl} % 表格单元格背景色支持

% ---------- 列表、代码与文本 ----------
\usepackage{enumitem} % 列表定制（间距、标签等）
\usepackage{soul} % 文本高亮（\hl）和删除线等
\usepackage{listings} % 源代码高亮与代码块样式
%========================================
% Emoji support using twemojis (Vector SVGs)
%========================================
\usepackage{twemojis} % 将 emoji 渲染为矢量图（twemoji SVG）
\usepackage{alltt} % 类似 verbatim 但允许 LaTeX 命令展开（用于示例输出）
\usepackage{alltt}

\usepackage{subfiles} % 支持将章节作为独立子文件编译
\usepackage[numbers,sort&compress]{natbib} % 参考文献样式与引用压缩

% ---------- 图题与浮动 ----------
\usepackage{float} % 强制浮动体位置（H）参数
\usepackage{caption} % 自定义图表题注样式
\usepackage{subcaption} % 子图题注环境

% ---------- 其他常用工具 ----------
\usepackage{changelog} % 变更日志支持（项目文档用途）
\usepackage{tcolorbox} % 彩色盒子用于提示、示例等
\tcbuselibrary{skins,breakable} % tcolorbox 的扩展库
\usepackage{etoolbox} % 编程型宏工具集合
\usepackage{siunitx} % 物理量与单位格式化

% ---------- TODO 注释 ----------
% 默认使用页边 TODO，尽量不影响正文排版。
\usepackage[textsize=scriptsize]{todonotes} % 文档内待办注释（边注）
\setlength{\marginparwidth}{2.1cm}

% ---- TikZ 外部化缓存 (首次编译生成 PDF，后续直接复用，大幅提速) ----
% 编译时必须开启 -shell-escape（VSCode 配置已设置）
% 如需临时关闭缓存（调试用），注释下面两行并取消 \tikzexternaldisable 的注释
% prefix 只能是项目根目录下的相对路径（不含 out/ 前缀），
% 因为 xelatex 的 -output-directory=out 会让 openout 自动在前面加 out/，
% 若 prefix 也写 out/ 则变成 out/out/tikz-cache/ —— 路径叠加导致写入失败。
% .md5 文件写到 out/tikz-cache/（kpathsea 会把 out/ 加入搜索路径，读取正常）；
% 外部化子进程（无 -output-directory）将 PDF 写到项目根 tikz-cache/，\includegraphics 可直接找到。
\tikzexternalize[prefix=tikz-cache/]
%\tikzexternaldisable  % 取消注释可关闭外部化，回退到内联编译（调试用）

% ---- 修复 tcolorbox (skins+breakable) 与 TikZ external 的兼容性冲突 ----
% tcolorbox skins 库内部会创建 tikzpicture 用于绘制边框/背景；
% TikZ external 全局 hook 会拦截这些内部 tikzpicture，在 breakable 跨页时
% 导致 collect@until@end@tikzpicture 状态错乱（Extra \else / extra } 等连锁错误）。
% 解决方案：在每个 tcolorbox / personalnote 环境的开头/结尾禁用/启用外部化。
\AtBeginEnvironment{tcolorbox}{\tikzexternaldisable}
\AtEndEnvironment{tcolorbox}{\tikzexternalenable}
\AtBeginEnvironment{personalnote}{\tikzexternaldisable}
\AtEndEnvironment{personalnote}{\tikzexternalenable}
\AtBeginEnvironment{taskbox}{\tikzexternaldisable}
\AtEndEnvironment{taskbox}{\tikzexternalenable}

\definecolor{ieeeBlue}{RGB}{0, 83, 159}
\definecolor{ieeeRed}{RGB}{198, 38, 46}
\definecolor{ieeeGreen}{RGB}{0, 114, 41}
\definecolor{ieeeOrange}{RGB}{235, 110, 31}
\definecolor{ieeeGray}{RGB}{100, 100, 100}
\definecolor{darkGray}{RGB}{66, 66, 66}
\definecolor{lightGray}{RGB}{189, 189, 189}

% TODO 全局样式与便捷命令
\setuptodonotes{
  color=ieeeOrange!15,
  bordercolor=ieeeOrange,
  linecolor=ieeeOrange
}
% todonotes 内部依赖 TikZ；开启 externalize 时需临时禁用外部化，避免编译失败。
\let\origTodo\todo
\RenewDocumentCommand{\todo}{ O{} m }{%
  \tikzexternaldisable
  \origTodo[#1]{#2}%
  \tikzexternalenable
}
\newcommand{\todoinline}[1]{\todo[inline,color=ieeeOrange!12]{#1}}

\usepackage{pgfplots} % 基于 TikZ 的函数与数据绘图
\pgfplotsset{compat=1.18}


% ctex 设置，确保各级标题缩进符合中文规范
\ctexset{
  paragraph/indent = 2em,
  subparagraph/indent = 2em,
}



% 链接相关宏包建议放在导言区靠后位置
\usepackage{hyperref}
\hypersetup{
  colorlinks=true,
  linkcolor=blue,
  urlcolor=blue,
  pdftitle={知识库文档},
  pdfauthor={葛良晨}
}

\lstdefinelanguage{riscv-asm}{
  morekeywords={add, addi, sub, sll, slli, srl, srli, sra, srai, and, andi, or, ori, xor, xori, slt, slti, sltu, sltiu,
    beq, bne, blt, bge, bltu, bgeu, jal, jalr, lui, auipc, lw, lh, lb, lhu, lbu, sw, sh, sb,
    csrr, csrw, csrrs, csrrc, csrrwi, csrrsi, csrrci, mret, ecall, ebreak, wfi, nop, ret, call, j, li, la, mv,
    .text, .data, .bss, .rodata, .section, .global, .globl, .align, .file, .type, .size, .word, .half, .byte, .asciz},
  morecomment=[l]{\#},
  morecomment=[l]{;},
  morestring=[b]",
  sensitive=true
}

\lstdefinelanguage{SystemVerilog}{
  morekeywords={
    always_comb, always_ff, always_latch, assert, assume, automatic, begin, end, case, endcase,
    class, endclass, clocking, endclocking, constraint, endconstraint, cover, covergroup, endgroup,
    default, disable, do, else, endfunction, endmodule, endpackage, endprogram, endtask, enum,
    export, extends, final, for, foreach, fork, join, join_any, join_none, function, generate,
    endgenerate, genvar, if, import, inside, interface, endinterface, localparam, logic, modport,
    module, new, package, program, protected, rand, randc, randcase, randsequence, ref, return,
    sequence, endsequence, struct, task, typedef, union, unique, unique0, virtual, void, wait, while,
    with
  },
  morecomment=[l]{//},
  morecomment=[s]{/*}{*/},
  morestring=[b]",
  sensitive=true
}

\lstset{
  basicstyle=\ttfamily\small,
  frame=single,
  breaklines=true,
  keywordstyle=\color{blue},
  commentstyle=\color{teal},
  stringstyle=\color{red!70!black}
}

% ---------- 自定义行内代码样式 ----------
\definecolor{codebg}{RGB}{230, 230, 230}
% 定义新命令 \code，避免直接修改 \texttt 导致系统崩溃
\newcommand{\code}[1]{\colorbox{codebg}{\ttfamily #1}}

% 保持 tcolorbox 的兼容性设置
\tcbset{shield externalize}





\title{RISC-V 技术文档 Demo}
\author{葛良晨}
\date{\today}

% 自定义个人总结环境
\newtcolorbox{personalnote}[1][]{
  colback=yellow!10,
  colframe=yellow!70!black,
  fonttitle=\bfseries,
  title=个人总结,
  breakable,
  enhanced,
  before upper={\setlength{\parindent}{2em}},
  #1
}

% 自定义待办事项环境
\newtcolorbox{taskbox}[1][]{
  colback=blue!5,
  colframe=blue!60!cyan!70!black,
  fonttitle=\bfseries,
  title=待办事项,
  breakable,
  enhanced,
  before upper={\setlength{\parindent}{2em}},
  #1
}

% 自定义妙妙想法彩框
\newtcolorbox{ideabox}[1][]{
  colback=purple!8!white,
  colframe=purple!60!black,
  fonttitle=\bfseries,
  title=✨ 妙妙想法,
  breakable,
  enhanced,
  before upper={\setlength{\parindent}{2em}},
  #1
}




\title{GLM-5.1 本地部署方案深度调研}
\author{葛良晨 by Claude Code}
\date{\today}

% 参考文献使用 natbib (numbers 格式)
% 由于本项目不使用 .bib 文件，采用行内 URL 引用方式

\begin{document}
\maketitle

% ------------------ 文档版本历史 ------------------
\section*{文档版本历史}
\begin{itemize}
    \item [2026-05-20] \textbf{重要勘误}：修正显存计算错误。经核查官方 vLLM recipe 页面（\href{https://github.com/vllm-project/recipes/blob/main/GLM/GLM5.md}{vllm-project/recipes}）明确标注目标硬件为"8xH200 (or H20) GPUs (141GB$\times$8)"，而非 H100。FP8 权重 93 GB/卡超过 H100 80GB 物理上限 16\%。全面重算各精度下的最低 GPU 配置，并补全所有信息来源引用。
\item [2026-05-19] 初始版本，完成 GLM-5.1 模型的本地部署方案全面调研。

\textbf{信息来源说明}：本文数据来源包括：(1) HuggingFace 模型页 \href{https://huggingface.co/zai-org/GLM-5.1}{zai-org/GLM-5.1}；(2) GitHub 官方仓库 \href{https://github.com/zai-org/GLM-5}{zai-org/GLM-5}；(3) vLLM Recipe 页面 \href{https://github.com/vllm-project/recipes/blob/main/GLM/GLM5.md}{vllm-project/recipes/GLM/GLM5.md}；(4) 模型 config.json \href{https://huggingface.co/zai-org/GLM-5.1/blob/main/config.json}{zai-org/GLM-5.1/config.json}；(5) 技术报告 \href{https://arxiv.org/abs/2602.15763}{arXiv:2602.15763}。总参数存在 744B（GitHub 官方）与 754B（HuggingFace 页面）两种说法，差异约 1.3\%，本文采用 GitHub 官方值 744B。激活参数 GitHub 官方表述为 40B。以下所有分析基于以上来源交叉验证。
\end{itemize}

% ======================================================================
% 第一章：GLM-5.1 模型概述
% ======================================================================

\section{GLM-5.1 模型概述}

\subsection{基本信息}

GLM-5.1 是智谱 AI（Z.ai / THUDM）于 2025 年中发布的新一代旗舰大语言模型，定位为"Agentic Engineering（智能体工程）"时代的基座模型。其核心定位是：\textbf{不是单纯的对话模型，而是面向长时间、多轮次、数千次工具调用的自主编程智能体}。

\begin{table}[htbp]
\centering
\caption{GLM-5.1 基本信息}
\begin{tabular}{p{3cm}p{9cm}}
\hline
\textbf{项目} & \textbf{详情} \\
\hline
开发者 & 智谱 AI (Z.ai / THUDM) \\
模型架构 & GLM MoE + DSA（DeepSeek Sparse Attention） \\
总参数量 & 744B（GitHub 官方）\textsuperscript{[1]}；754B（HuggingFace 页面显示）\textsuperscript{[2]} \\
激活参数量 & 40B（GitHub 官方），每次推理仅激活约 5.4\% 的参数 \\
精度格式 & BF16（原始）、FP8（官方优化版） \\
上下文窗口 & 200K tokens（config.json: max\_position\_embeddings=202752）\textsuperscript{[4]} \\
架构细节 & 78 层，75 层 MoE（每层 256 路由专家 + 1 共享专家），3 层 Dense FFN \\
MLA 注意力 & KV LoRA rank=512，Q LoRA rank=2048，压缩比约 93\% \\
开源协议 & MIT\textsuperscript{[2]} / Apache 2.0\textsuperscript{[1]} \\
技术报告 & \href{https://arxiv.org/abs/2602.15763}{arXiv:2602.15763} \textsuperscript{[5]} \\
模型页 & \href{https://huggingface.co/zai-org/GLM-5.1}{huggingface.co/zai-org/GLM-5.1} \textsuperscript{[2]} \\
官方仓库 & \href{https://github.com/zai-org/GLM-5}{github.com/zai-org/GLM-5} \textsuperscript{[1]} \\
\hline
\end{tabular}
\end{table}

\subsection{架构特点：MLA + MoE + DSA}

GLM-5.1 的架构标签为 \texttt{glm\_moe\_dsa}，融合了三项核心技术：

\textbf{MLA（Multi-head Latent Attention，多头潜在注意力）}：源自 DeepSeek 系列（V2/V3）的注意力机制。传统 MHA 为每个 token 存储完整的 KV 矩阵；MLA 将 K 和 V 联合压缩到低维潜在空间（KV LoRA rank=512），仅在计算时动态解压。这一设计将 KV Cache 压缩约 93\%，是 GLM-5.1 在 200K 长上下文下仍能控制显存的关键。

\textbf{MoE（Mixture of Experts，混合专家）}：78 层中 75 层采用 MoE，每层包含 1 个共享专家（始终激活）+ 256 个路由专家。每次推理仅激活 8 个路由专家 + 1 个共享专家，使得计算量仅与约 40B 参数的稠密模型相当，而知识容量达到 744B 级别。路由采用 sigmoid 门控 + token-choice 策略，无辅助损失。

\textbf{DSA（DeepSeek Sparse Attention，深度稀疏注意力）}：每层配备一个 Indexer（32 头、128 维），从 200K 上下文中智能选择 top-2048 个最相关的 KV 对。这避免了在超长上下文中对所有 token 做全量 Attention 的 $O(L^2)$ 复杂度。

三者协同工作的效果：MLA 压缩 KV Cache 体积，DSA 稀疏化 Attention 计算范围，MoE 将知识容量与推理成本解耦。这就是 GLM-5.1 能以"40B 级别的推理成本"承载"744B 级别的知识容量"的底层原因。

\subsection{性能定位}

GLM-5.1 在多个权威基准上达到了开源模型的最高水平，甚至在部分指标上超越了闭源商业模型。

\begin{table}[htbp]
\centering
\caption{GLM-5.1 核心基准测试成绩}
\begin{tabular}{p{4.5cm}p{2cm}p{5.5cm}}
\hline
\textbf{基准测试} & \textbf{得分} & \textbf{排名/备注} \\
\hline
SWE-Bench Pro（软件工程） & 58.4 & \textbf{\#1 SOTA}，超越 Claude Opus 4.6 \\
CyberGym（网络安全） & 68.7 & \textbf{\#1 SOTA} \\
BrowseComp（网页理解） & 68.0 & \textbf{\#1 SOTA} \\
NL2Repo（自然语言到仓库） & 42.7 & \#2，仅次于 Claude Opus 4.6 (49.8) \\
Terminal-Bench 2.0 & 63.5 & 终端环境编程能力 \\
AIME 2026（数学竞赛） & 95.3 & 数学推理 \\
HMMT Feb 2026（数学竞赛）& 82.6 & 数学推理 \\
GPQA-Diamond（研究生知识）& 86.2 & 知识推理 \\
HLE（难例评测） & 31.0 & 带工具提升至 52.3 \\
MCP-Atlas（工具调用） & 71.8 & 多工具协同 \\
\hline
\end{tabular}
\end{table}

\textbf{关键差异化能力}：GLM-5.1 的设计重心并非单轮对话质量，而是"the longer it runs, the better the result"——在数百轮迭代、数千次工具调用的长时间自主编程任务中，模型不会像前代那样过早陷入局部最优或"迷失方向"。这一特性使其在 SWE-Bench Pro 等需要持续优化的复杂工程任务上取得突破。

\subsection{与竞品模型的参数对比}

\begin{table}[htbp]
\centering
\caption{GLM-5.1 与主流开源 MoE 模型参数对比}
\begin{tabular}{p{2.5cm}p{1.8cm}p{1.8cm}p{2cm}p{2.5cm}}
\hline
\textbf{模型} & \textbf{总参数} & \textbf{激活参数} & \textbf{激活比} & \textbf{上下文} \\
\hline
GLM-5.1 & 744B & 40B & 5.4\% & 128K \\
GLM-5 & 744B & 40B & 5.4\% & 128K \\
GLM-4.5 & 355B & 32B & 9.0\% & 128K \\
DeepSeek-V3 & 671B & 37B & 5.5\% & 128K \\
Qwen3-235B & 235B & 22B & 9.4\% & 128K \\
Mixtral 8$\times$22B & 141B & 39B & 27.7\% & 64K \\
\hline
\end{tabular}
\end{table}

GLM-5.1 的激活比（5.4\%）在同类 MoE 模型中处于极低水平，意味着相同推理成本下可以承载更大的知识容量。但这也意味着\textbf{模型权重的静态存储需求极大}——744B 参数即使不全部激活，也必须全部加载到显存/内存中。

% ======================================================================
% 第二章：硬件需求分析
% ======================================================================

\section{硬件需求分析}

\subsection{显存需求的理论计算}

GLM-5.1 作为 744B 参数的 MoE 模型，其显存需求比同级别稠密模型更为复杂。以下是完整的显存占用分解。

\subsubsection{权重部分的显存占用}

\begin{table}[htbp]
\centering
\caption{不同精度下 GLM-5.1 权重的显存占用}
\begin{tabular}{p{2cm}p{3cm}p{3.5cm}p{3.5cm}}
\hline
\textbf{精度} & \textbf{每参数字节} & \textbf{纯权重占用} & \textbf{说明} \\
\hline
BF16 & 2 bytes & $\approx$1,488 GB & 原始精度，不可单机承载 \\
FP8 & 1 byte & $\approx$744 GB & 官方优化版，需 8$\times$H200 \\
INT8 & 1 byte & $\approx$744 GB & 与 FP8 等同 \\
INT4 & 0.5 bytes & $\approx$372 GB & GGUF Q4\_K\_M 典型值 \\
NVFP4 & $\sim$0.59 bytes & $\approx$437 GB & NVIDIA Blackwell FP4 专用格式 \\
MXFP4 & $\sim$0.55 bytes & $\approx$408 GB & AMD MXFP4 专用格式 \\
INT2 & 0.25 bytes & $\approx$186 GB & 极低精度，质量损失严重 \\
\hline
\end{tabular}
\end{table}

\subsubsection{完整显存账单（推理时）——关键计算}

推理时的总显存占用远不止权重本身。以一个实际的部署场景为例：

\begin{equation}
\begin{aligned}
\text{VRAM}_{\text{total}} = &\ \text{权重} + \text{KV Cache} + \text{激活值} + \text{框架开销} \\
                            &\ + \text{MoE 专家路由缓冲区} + \text{通信缓冲区}
\end{aligned}
\end{equation}

\textbf{以官方 vLLM recipe 的实际配置（8$\times$H200, FP8, TP=8）为例}\textsuperscript{[3]}：

\begin{itemize}
    \item \textbf{权重}：744 GB（FP8），Tensor Parallelism 均分到 8 张 GPU，每张约 $\mathbf{93\ GB}$
    \item \textbf{KV Cache}（200K 上下文时的峰值）：得益于 MLA（KV LoRA rank=512）和 DSA 稀疏注意力（每查询仅保留 top-2048 个 KV），KV Cache 远小于标准 Attention。MLA 压缩后每 token 的 KV 约 2 KB，200K 上下文单请求约 0.4 GB/GPU；但并发场景下需按请求数累乘
    \item \textbf{MoE 专家路由缓冲区}：约 4--6 GB/GPU（路由过程中的临时激活值）
    \item \textbf{框架开销}（CUDA context、vLLM 内核缓存）：约 2--3 GB/GPU
    \item \textbf{通信缓冲区}（TP all-reduce 缓冲）：约 1--2 GB/GPU
\end{itemize}

\textbf{单卡峰值显存需求（FP8, 200K 上下文, 单请求）}：
$93 + 0.4 + 6 + 3 + 2 \approx \mathbf{104\ GB}$

\textbf{关键判断}：
\begin{itemize}
    \item \textbf{H200 141 GB}：104 GB $<$ 141 GB，剩余 $\sim$37 GB 可承载并发 KV Cache。\textbf{可行。}这正是官方 vLLM recipe 的目标配置。
    \item \textbf{H100/A100 80 GB}：104 GB $>$ 80 GB。\textbf{即使不计 KV Cache，93 GB 纯权重已经超过 80 GB 物理上限。不可行。}
\end{itemize}

\textbf{为什么 8$\times$H100 跑不了 FP8——用数字说话}（来源\textsuperscript{[3]}确认 H200 为目标硬件，结合 H100 物理显存计算）：

FP8 权重 744 GB，纯 Tensor Parallelism（TP=8）下每卡承载 $744/8 = 93\text{ GB}$。H100 物理显存为 80 GB。\textbf{仅权重一项即超出单卡容量 16\%}。这不是"紧张"或"临界"——是物理上无法完整加载模型。这就是官方 vLLM recipe 使用 H200（141 GB）而非 H100 的根本原因。

\textbf{如果坚持用 H100 80GB 部署，必须满足以下条件之一}：
\begin{enumerate}
    \item \textbf{增加 TP 规模}：TP $\geq$ 12（每卡 $744/12 = 62$ GB，剩余 18 GB 勉强够短上下文单请求）
    \item \textbf{使用流水线并行（PP）}：PP=2 + TP=8，共 16 卡，每层权重分布到更少的 GPU
    \item \textbf{使用更低精度}：Q4\_K\_M（$\sim$465 GB），TP=8 每卡 $\sim$58 GB，可在 8$\times$A100 上运行
    \item \textbf{引入 CPU offload}：将部分 MoE 专家层卸载到系统内存（见方案 C）
\end{enumerate}

\subsection{GPU 选型分析}

\subsubsection{候选 GPU 规格对比}

\begin{table}[htbp]
\centering
\caption{适用于 GLM-5.1 部署的 GPU 规格对比}
\begin{tabular}{p{2.2cm}p{1.6cm}p{2cm}p{1.6cm}p{1.8cm}p{1.6cm}}
\hline
\textbf{GPU} & \textbf{显存} & \textbf{带宽} & \textbf{FP16算力} & \textbf{互联} & \textbf{单卡价格} \\
\hline
H200 SXM & 141 GB & 4.8 TB/s & 990 TFLOPS & NVLink 900GB/s & \textasciitilde\$35K \\
H100 SXM & 80 GB & 3.35 TB/s & 990 TFLOPS & NVLink 900GB/s & \textasciitilde\$28K \\
H100 PCIe & 80 GB & 2.0 TB/s & 756 TFLOPS & 无 NVLink & \textasciitilde\$24K \\
A100 SXM & 80 GB & 2.0 TB/s & 312 TFLOPS & NVLink 600GB/s & \textasciitilde\$15K \\
A100 PCIe & 80 GB & 1.9 TB/s & 312 TFLOPS & 无 NVLink & \textasciitilde\$12K \\
A6000 Ada & 48 GB & 0.96 TB/s & 91 TFLOPS & 无 & \textasciitilde\$6.8K \\
L40S & 48 GB & 0.86 TB/s & 91 TFLOPS & 无 & \textasciitilde\$8K \\
RTX 6000 Ada & 48 GB & 0.96 TB/s & 91 TFLOPS & 无 & \textasciitilde\$6.8K \\
RTX 4090 & 24 GB & 1.0 TB/s & 83 TFLOPS & 无 & \textasciitilde\$1.6K \\
\hline
\end{tabular}
\end{table}

\subsubsection{显存容量约束分析——基于物理上限的正确计算}

GLM-5.1 的部署核心约束是\textbf{单 GPU 显存容量必须 $\geq$ 权重分片 + 运行时开销}，而非简单的"总显存大于总权重"。

\begin{table}[htbp]
\centering
\caption{不同 GPU 配置的 GLM-5.1 部署可行性（基于物理上限）}
\begin{tabular}{p{2.2cm}p{1.4cm}p{2cm}p{2cm}p{4cm}}
\hline
\textbf{GPU} & \textbf{数量} & \textbf{单卡显存} & \textbf{每卡权重} & \textbf{判断} \\
\hline
\multicolumn{5}{c}{\textbf{FP8 精度（权重共 744 GB）}} \\
\hline
H200 SXM & 8 & 141 GB & 93 GB & \textbf{可行（官方配置）}：48 GB 余量 \\
H100 SXM & 8 & 80 GB & 93 GB & \textbf{不可行}：权重超限 16\% \\
H100 SXM & 12 & 80 GB & 62 GB & \textbf{临界可行}：18 GB 余量，仅短上下文 \\
H100 SXM & 16 & 80 GB & 46.5 GB & \textbf{可行}：33.5 GB 余量，舒适 \\
A100 SXM & 12 & 80 GB & 62 GB & \textbf{临界可行}：带宽低于 H100，速度受限 \\
A100 SXM & 16 & 80 GB & 46.5 GB & \textbf{可行}：推荐 A100 最低配置 \\
A6000 Ada & 16 & 48 GB & 46.5 GB & \textbf{不可行}：余量仅 1.5 GB，无 KV Cache 空间 \\
\hline
\multicolumn{5}{c}{\textbf{BF16 精度（权重共 1,488 GB）}} \\
\hline
H200 SXM & 8 & 141 GB & 186 GB & \textbf{不可行}：权重超限 32\% \\
H200 SXM & 16 & 141 GB & 93 GB & \textbf{可行}：48 GB 余量，推荐 BF16 配置 \\
H200 SXM & 12 & 141 GB & 124 GB & \textbf{临界可行}：17 GB 余量 \\
H100 SXM & 24 & 80 GB & 62 GB & \textbf{可行}：18 GB 余量，需 NVLink 全互联 \\
\hline
\multicolumn{5}{c}{\textbf{Q4\_K\_M GGUF（权重共 $\sim$465 GB）}} \\
\hline
A100 SXM & 8 & 80 GB & 58 GB & \textbf{可行}：22 GB 余量，短上下文 OK \\
A100 SXM & 12 & 80 GB & 39 GB & \textbf{舒适}：41 GB 余量，支持长上下文 \\
A6000 Ada & 12 & 48 GB & 39 GB & \textbf{临界可行}：9 GB 余量 \\
\hline
\end{tabular}
\end{table}

\textbf{核心结论（修正后）}：
\begin{itemize}
    \item \textbf{FP8 官方最低配置}：\textbf{8$\times$H200 (141 GB)}，不是 8$\times$H100
    \item \textbf{FP8 + H100/A100 最低配置}：至少 \textbf{12 张}（80 GB），推荐 16 张
    \item \textbf{BF16 最低配置}：\textbf{12--16$\times$H200}，或 \textbf{24$\times$H100}
    \item \textbf{Q4\_K\_M GGUF 最低配置}：\textbf{8$\times$A100 (80 GB)}，这是当前自建部署的性价比甜点
    \item \textbf{48 GB 级 GPU 无法承载 FP8}：即使 16 张也几乎没有 KV Cache 余量
\end{itemize}

\subsection{量化与显存的精确对应关系}

量化是降低 GLM-5.1 部署门槛的核心手段。以下基于社区已有的 GGUF 量化数据（以 unsloth/GLM-5.1-GGUF 为例）：

\begin{table}[htbp]
\centering
\caption{GLM-5.1 GGUF 量化变体的文件大小与显存需求（修正版）}
\begin{tabular}{p{2.5cm}p{2.2cm}p{2cm}p{1.8cm}p{3.2cm}}
\hline
\textbf{量化级别} & \textbf{文件总大小} & \textbf{TP=8 每卡} & \textbf{单卡需求} & \textbf{推荐 GPU 配置} \\
\hline
BF16（原始） & $\sim$1,488 GB & 186 GB & 200+ GB & 16$\times$H200 或 24$\times$H100 \\
FP8（官方） & $\sim$744 GB & 93 GB & 115+ GB & \textbf{8$\times$H200}（官方 recipe）\\
Q8\_0 & $\sim$744 GB & 93 GB & 115+ GB & 8$\times$H200（GGUF 无 TP 加速）\\
Q6\_K & $\sim$558 GB & 70 GB & 90+ GB & 8$\times$H200 或 12$\times$H100 \\
Q5\_K\_M & $\sim$465 GB & 58 GB & 78+ GB & 8$\times$A100（80GB）\\
Q4\_K\_M & $\sim$465 GB & 58 GB & 78+ GB & \textbf{8$\times$A100}（性价比甜点）\\
Q3\_K\_M & $\sim$372 GB & 47 GB & 65+ GB & 8$\times$A100 或 12$\times$A6000 \\
NVFP4 & $\sim$437 GB & 55 GB & 75+ GB & 8$\times$H100/B100 \\
MXFP4 & $\sim$408 GB & 51 GB & 70+ GB & 8$\times$MI300X \\
IQ4\_XS & $\sim$372 GB & 47 GB & 65+ GB & 12$\times$A6000 或 8$\times$A100 \\
IQ3\_XXS & $\sim$279 GB & 35 GB & 50+ GB & 8$\times$A6000（48GB）\\
IQ2\_XXS & $\sim$186 GB & 23 GB & 40+ GB & 8$\times$A6000 或 CPU offload \\
\hline
\end{tabular}
\end{table}

\textbf{注意}：以上"单卡需求"包含了 KV Cache、框架开销和约 15\% 的安全余量。llama.cpp 的层级拆分（layer-wise splitting）机制与 vLLM 的 TP 不同——每张 GPU 加载若干完整层而非每层的 1/N——但在 GGUF 格式下，多 GPU 加载等价于将所有层均匀分布到多张 GPU，因此单卡需求的计算方式类似。使用 CPU offload 可进一步降低 GPU 显存需求，但会显著降低吞吐量。

% ======================================================================
% 第三章：部署方案详析
% ======================================================================

\section{部署方案详析}

\subsection{方案 A：企业级全精度部署（BF16 + 16$\times$H200）}

这是性能最完整、成本也最高的方案。适用于对模型质量有极致要求的企业级生产环境。

\textbf{硬件配置}：
\begin{itemize}
    \item 16$\times$ NVIDIA H200 SXM（141 GB 每卡），总显存 2,256 GB
    \item 或 24$\times$ H100 SXM（80 GB 每卡），总显存 1,920 GB
    \item NVLink 全互联 + NVSwitch 拓扑（跨 16/24 卡需要多节点 InfiniBand 互联）
    \item BF16 权重 1,488 GB，TP=16 每卡 93 GB，剩余 48 GB 给 KV Cache 和运行时开销
    \item \textbf{注意}：8$\times$H200 无法运行 BF16——1,488/8 = 186 GB/卡 $>$ 141 GB 物理上限
\end{itemize}

\textbf{部署命令（vLLM, 16$\times$H200, 参考官方 recipe 调整）}\textsuperscript{[3]}：
\begin{lstlisting}[language=bash]
vllm serve zai-org/GLM-5.1 \
    --tensor-parallel-size 16 \
    --speculative-config.method mtp \
    --speculative-config.num_speculative_tokens 3 \
    --tool-call-parser glm47 \
    --reasoning-parser glm45 \
    --enable-auto-tool-choice \
    --chat-template-content-format=string \
    --served-model-name glm-5.1-bf16
\end{lstlisting}

\textbf{预期性能}：
\begin{itemize}
    \item 单请求 decode 速度：50--80 tokens/s（16$\times$H200, BF16）
    \item 并发能力：10--20 个并发请求（得益于大容量 KV Cache 空间）
    \item Prefill 速度：128K token Prompt $\approx$ 3--5 秒
\end{itemize}

\textbf{优势}：全精度无质量损失，NVLink 保证 TP 通信效率，大量 KV Cache 余量支持高并发。\newline
\textbf{劣势}：硬件成本极高（16$\times$H200 约 \$56 万），功耗巨大（约 11 kW），需要多节点数据中心级部署。

\subsection{方案 B：官方 FP8 量化部署（8$\times$H200 或 12--16$\times$H100）}

这是官方 vLLM recipe 的标准部署方案，在精度与成本之间取得了工程上的最优平衡。官方提供了预量化的 FP8 模型（\texttt{zai-org/GLM-5.1-FP8}）。

\textbf{⚠ 重要勘误}：本文初版错误地认为 8$\times$H100 (80 GB) 可以承载 FP8 部署。错误来源：GitHub README\textsuperscript{[1]} 的部署命令给出了 \texttt{--tensor-parallel-size 8} 和 \texttt{--gpu-memory-utilization 0.85}，但\textbf{未指定 GPU 型号}；HuggingFace 页面\textsuperscript{[2]} 的命令完全省略了 TP 配置；仅 vLLM recipe 页面\textsuperscript{[3]} 明确标注 "Serving FP8 Model on 8xH200 (or H20) GPUs (141GB$\times$8)"。三份官方文档的详细程度差异是初版判断错误的直接原因。实际上 FP8 权重 93 GB/卡 超过 H100 80 GB 物理上限 16\%。官方 vLLM recipe 的目标硬件是 \textbf{H200 (141 GB)}。

\textbf{硬件配置——三个选项}：

\textbf{选项 B-1：官方推荐（8$\times$H200）}
\begin{itemize}
    \item 8$\times$ NVIDIA H200 SXM（141 GB 每卡），总显存 1,128 GB
    \item FP8 权重 744 GB，TP=8 每卡 93 GB，剩余 48 GB 给 KV Cache 和开销
    \item \textbf{这是官方 vLLM recipe 的精确配置}
\end{itemize}

\textbf{选项 B-2：H100 扩容方案（16$\times$H100）}
\begin{itemize}
    \item 16$\times$ NVIDIA H100 SXM（80 GB 每卡），总显存 1,280 GB
    \item FP8 权重 744 GB，TP=16 每卡 46.5 GB，剩余 33.5 GB 给 KV Cache 和开销
    \item NVLink 全互联（900 GB/s），但跨 16 卡需要 InfiniBand 跨节点通信
\end{itemize}

\textbf{选项 B-3：H100 最低方案（12$\times$H100）}
\begin{itemize}
    \item 12$\times$ NVIDIA H100 SXM（80 GB 每卡），总显存 960 GB
    \item FP8 权重 744 GB，TP=12 每卡 62 GB，剩余 18 GB 给 KV Cache 和开销
    \item \textbf{临界配置}：KV Cache 空间紧张，仅适合短上下文（$\leq$32K）或单请求场景
\end{itemize}

\textbf{部署命令（SGLang, 8$\times$H200, 来自 GitHub 官方 README）}\textsuperscript{[1]}：
\begin{lstlisting}[language=bash]
sglang serve \
    --model-path zai-org/GLM-5.1-FP8 \
    --tp-size 8 \
    --tool-call-parser glm47 \
    --reasoning-parser glm45 \
    --speculative-algorithm EAGLE \
    --speculative-num-steps 3 \
    --speculative-eagle-topk 1 \
    --speculative-num-draft-tokens 4 \
    --mem-fraction-static 0.85 \
    --served-model-name glm-5.1-fp8 \
    --port 8000 \
    --host 0.0.0.0
\end{lstlisting}

\textbf{vLLM 官方 recipe 命令（8$\times$H200 或 H20，来自 vLLM recipe 页面）}\textsuperscript{[3]}：
\begin{lstlisting}[language=bash]
vllm serve zai-org/GLM-5.1-FP8 \
    --tensor-parallel-size 8 \
    --speculative-config.method mtp \
    --speculative-config.num_speculative_tokens 3 \
    --tool-call-parser glm47 \
    --reasoning-parser glm45 \
    --enable-auto-tool-choice \
    --chat-template-content-format=string \
    --served-model-name glm-5.1-fp8
\end{lstlisting}

\textbf{预期性能}：
\begin{itemize}
    \item Decode 速度：40--60 tokens/s（8$\times$H200），30--45 tokens/s（16$\times$H100）
    \item 并发能力：8--15 个并发请求（8$\times$H200），5--10（16$\times$H100）
    \item Prefill 速度：128K token $\approx$ 4--6 秒
\end{itemize}

\textbf{优势}：官方直接支持，FP8 质量损失可忽略（$<0.5\%$ 困惑度差异），社区最成熟的部署路径。\newline
\textbf{劣势}：至少 8$\times$H200（\$28 万+），或 12--16$\times$H100（\$34--45 万），硬件门槛极高。

\subsection{方案 C：GGUF 量化方案（Q4\_K\_M + llama.cpp/Ollama）}

对于没有 8 卡企业级 GPU 的团队，GGUF 量化是降低硬件门槛的主要途径。

\textbf{硬件配置选项}：

\textbf{C-1：8$\times$A100 全 GPU 模式}
\begin{itemize}
    \item GGUF Q4\_K\_M 文件大小 $\sim$465 GB，8$\times$80 GB 全 GPU 加载
    \item llama.cpp CUDA backend，tensor-parallel 支持
    \item 预期速度：15--25 tokens/s
\end{itemize}

\textbf{C-2：4$\times$A100/A6000 + CPU offload}
\begin{itemize}
    \item 4 张 GPU 承载最活跃的 MoE 专家层（约 60\% 权重）
    \item 剩余 40\% 权重 offload 到系统内存（DDR5 512+ GB）
    \item 预期速度：5--10 tokens/s（受 CPU-GPU 数据传输瓶颈）
\end{itemize}

\textbf{C-3：纯 CPU 模式（仅概念验证）}
\begin{itemize}
    \item 需要 $\sim$500 GB 以上系统内存（如双路 EPYC + 768 GB DDR5）
    \item Q4\_K\_M 纯 CPU 推理
    \item 预期速度：0.5--2 tokens/s（仅可用于验证模型是否加载成功，不可用于实际使用）
\end{itemize}

\textbf{部署命令（llama.cpp with CUDA）}：
\begin{lstlisting}[language=bash]
# 下载 Q4_K_M GGUF 文件（以 unsloth 为例）
# 来自 https://huggingface.co/unsloth/GLM-5.1-GGUF

# 多 GPU 推理（4卡示例，其余 offload 到 CPU）
./llama-cli \
    -m GLM-5.1-UD-Q4_K_M-00001-of-00011.gguf \
    -ngl 60 \
    -c 32768 \
    -fa 1 \
    --tensor-split 24,24,24,24
\end{lstlisting}

\textbf{说明}：\texttt{-ngl 60} 表示将前 60 层（MoE 中激活频率最高的部分）加载到 GPU，其余层留在 CPU。\texttt{--tensor-split} 指定每张 GPU 的显存配额（单位 GB）。对于 GLM-5.1 这样的 MoE 模型，GGUF 的 CPU offload 性能取决于 MoE 路由的频率分布——如果大部分 token 都路由到 GPU 驻留的专家，则性能接近纯 GPU；否则会频繁触发 CPU-GPU 数据搬运。

\subsection{方案 D：NVFP4 / MXFP4 新精度格式}

这是 NVIDIA 和 AMD 分别推出的 4-bit 浮点格式，相比传统 INT4 量化有更好的动态范围和精度保持。

\textbf{D-1：NVIDIA NVFP4（Blackwell 架构专用）}

NVFP4 是 NVIDIA Blackwell（B100/B200）架构引入的原生 4-bit 浮点格式。与 INT4 不同，NVFP4 保留了浮点的指数位，在极低精度下仍能覆盖较大的数值范围。
\begin{itemize}
    \item 模型大小：$\sim$437 GB（社区已有 lukealonso/GLM-5.1-NVFP4）
    \item 需 Blackwell 架构 GPU（B100/B200），利用硬件原生 FP4 Tensor Core
    \item 8$\times$B100 或 4$\times$B200（B200 显存更大）
    \item 理论吞吐量可达 FP8 的 2$\times$（受益于 4-bit 内存带宽减半）
\end{itemize}

\textbf{D-2：AMD MXFP4（MI300X/MI350 专用）}

AMD 的 MXFP4 是面向 ROCm 生态的 4-bit 浮点格式。
\begin{itemize}
    \item 模型大小：$\sim$408 GB（社区已有 amd/GLM-5.1-MXFP4）
    \item 需 AMD Instinct MI300X（192 GB）或更新型号
    \item 4$\times$MI300X（共 768 GB）即可承载，且有较大 KV Cache 余量
\end{itemize}

\subsection{方案 E：Apple Silicon 方案（MLX 量化）}

Apple Silicon 的统一内存架构使其在超大模型推理上有独特优势，但 GLM-5.1 的体量带来严峻挑战。

\begin{table}[htbp]
\centering
\caption{Apple Silicon 运行 GLM-5.1 的可行性分析}
\begin{tabular}{p{3cm}p{2cm}p{2.2cm}p{2.5cm}p{2cm}}
\hline
\textbf{设备} & \textbf{最大统一内存} & \textbf{内存带宽} & \textbf{可跑量化} & \textbf{预期速度} \\
\hline
Mac Studio M2 Ultra & 192 GB & 800 GB/s & MLX 2.5-bit & 2--4 tok/s \\
Mac Pro M2 Ultra & 192 GB & 800 GB/s & MLX 2.5-bit & 2--4 tok/s \\
Mac Studio M4 Max & 128 GB & 546 GB/s & \textbf{不可行} & --- \\
MacBook Pro M4 Max & 128 GB & 546 GB/s & \textbf{不可行} & --- \\
\hline
\end{tabular}
\end{table}

\textbf{关键限制}：
\begin{itemize}
    \item 即使用 MLX 2.5-bit 量化（$\sim$230 GB），也需要 192 GB 统一内存的设备，且几乎没有 KV Cache 余量（仅能运行极短上下文）
    \item Mac Pro M2 Ultra 是唯一"勉强可跑"的 Apple 设备，但速度（2--4 tokens/s）无法满足实际使用
    \item Apple Silicon 不支持 CUDA，MLX 对 GLM-5.1 的 MoE 架构优化尚不成熟
    \item \textbf{结论：Apple Silicon 不适合 GLM-5.1 的本地部署，不推荐此方案}
\end{itemize}

\subsection{方案 F：云服务方案}

对于缺少硬件预算的团队，云服务是最快上手的方式。

\textbf{F-1：Together AI API（推荐）}

Together AI 是 GLM-5.1 的官方推理服务商之一。
\begin{itemize}
    \item 提供 OpenAI 兼容的 Chat Completions API
    \item 无需管理基础设施，零部署成本
    \item 价格：按 token 计费（具体价格以官方页面为准，参考同类 700B+ 模型约 \$3--8/1M input tokens, \$8--15/1M output tokens）
\end{itemize}

\textbf{F-2：HuggingFace Inference Endpoint}

通过 HuggingFace 的推理端点直接部署 GLM-5.1。
\begin{itemize}
    \item 按 GPU 实例小时计费
    \item 以 8$\times$H100 实例为例，约 \$25--40/小时
    \item 适合间歇性使用、批量评估、或作为 PoC 阶段的过渡方案
\end{itemize}

\textbf{F-3：算力租赁（RunPod / Vast.ai / AutoDL）}

\begin{table}[htbp]
\centering
\caption{主流算力租赁平台对比（2026 年参考价格）}
\begin{tabular}{p{2.5cm}p{2.5cm}p{2.5cm}p{4cm}}
\hline
\textbf{平台} & \textbf{H100 80GB/时} & \textbf{A100 80GB/时} & \textbf{备注} \\
\hline
RunPod & \$2.49--3.49 & \$1.49--2.29 & 按需和竞价实例 \\
Vast.ai & \$1.80--3.00 & \$1.00--1.80 & 社区托管，价格波动大 \\
AutoDL（国内） & \$2.00--3.50 & \$1.20--2.00 & 国内部署首选，中文支持好 \\
Lambda Labs & \$2.49--2.99 & \$1.49--1.99 & 北美地区 \\
\hline
\end{tabular}
\end{table}

\textbf{F-4：HuggingChat（免费）}

GLM-5.1 已直接部署于 HuggingChat（huggingface.co/chat），可免费使用，适合快速体验模型能力，无需任何部署。

% ======================================================================
% 第四章：推理框架选型
% ======================================================================

\section{推理框架选型}

\subsection{框架能力矩阵}

\begin{table}[htbp]
\centering
\caption{GLM-5.1 支持的推理框架对比}
\begin{tabular}{p{2.5cm}p{2cm}p{3cm}p{2.8cm}p{1.8cm}}
\hline
\textbf{框架} & \textbf{最低版本} & \textbf{核心优势} & \textbf{适用场景} & \textbf{成熟度} \\
\hline
vLLM & v0.20.2\textsuperscript{[1]} & 生产级吞吐优化，PagedAttention，MTP推测解码 & 高并发在线服务 & 成熟 \\
SGLang & v0.5.11\textsuperscript{[1]} & RadixAttention，EAGLE推测解码 & 低延迟服务 & 成熟 \\
llama.cpp & 最新版 & GGUF 量化，CPU-GPU 混合 & 个人/小团队 & 成熟 \\
Ollama & 最新版 & 一键部署，自动下载 & 个人快速体验 & 成熟 \\
KTransformers & v0.5.3+ & 异构计算（GPU+CPU+NPU） & 降低硬件门槛 & 实验性 \\
xLLM & v0.8.0+ & 华为昇腾硬件 & 国产化替代 & 发展中 \\
Transformers & v0.5.3+ & HuggingFace 原生 & 研究/原型验证 & 基准 \\
\hline
\end{tabular}
\end{table}

\subsection{vLLM（官方首推）}

vLLM 是目前生产环境部署 GLM-5.1 的首选框架。其核心优势在于：

\textbf{PagedAttention}：将 KV Cache 按页（page）管理，类似操作系统的虚拟内存。这避免了传统框架中的显存碎片化问题，允许更高效的 batch 调度，将显存利用率从传统框架的 30--50\% 提升至 80--90\%。

\textbf{Continuous Batching}：动态地将新请求插入正在进行的 batch 中，无需等待整个 batch 完成，大幅提升吞吐量。

\textbf{MTP（Multi-Token Prediction）}：GLM-5.1 在 vLLM 中支持 MTP 推测解码，每次可预测 3 个 token，将 decode 吞吐量提升 30--50\%。

\textbf{生产特性}：OpenAI 兼容 API、Prometheus 监控指标、量化支持（FP8/INT8/AWQ）、前缀缓存（Prefix Caching）。

\subsection{SGLang}

SGLang 是来自 LMSys 的新一代推理框架，在延迟敏感场景下表现优异。

\textbf{RadixAttention}：基于前缀树（Radix Tree）的 KV Cache 复用。当多个请求共享相同的 system prompt 或文档前缀时，RadixAttention 自动共享 KV Cache，避免重复计算。对于 Agent 场景（每次工具调用后重新发送完整历史），这一特性尤为关键。

\textbf{EAGLE 推测解码}：比 MTP 更激进的推测解码策略，GLM-5.1 在 SGLang 上使用 EAGLE（3 步，topk=1，4 个草稿 token），在兼容性上优于 MTP。

\textbf{Docker 一键部署}：SGLang 提供了完整的 Docker 镜像，开箱即用。

\subsection{llama.cpp / Ollama}

作为本地 LLM 推理的基石，llama.cpp 在 GLM-5.1 的 GGUF 量化方案中不可替代。

\textbf{关键能力}：
\begin{itemize}
    \item \textbf{CPU-GPU 混合推理}：可将部分层 offload 到系统内存，大幅降低 GPU 需求
    \item \textbf{跨平台}：支持 CUDA、Metal（Apple Silicon）、Vulkan、ROCm
    \item \textbf{任意量化级别}：从 IQ1\_M（$\sim$1.7 bpw）到 Q8\_0（8 bpw）的完整光谱
    \item \textbf{上下文压缩}：支持 Flash Attention，降低长上下文的 KV Cache 开销
\end{itemize}

\textbf{Ollama} 在 llama.cpp 之上封装了类 Docker 的使用体验，但对于 GLM-5.1 这样的超大规模模型，Ollama 的默认配置可能不足，建议直接使用 llama.cpp 进行精细调参。

\subsection{KTransformers（异构计算方案）}

KTransformers 由 kvcache-ai 团队开发，专注于利用异构计算资源降低超大模型的部署门槛。

\textbf{核心思路}：将 MoE 模型中激活频率低的专家层 offload 到 CPU 甚至是 NPU，仅在 GPU 上保留高频激活的专家和 Attention 层。对于 GLM-5.1（激活比仅 5.4\%），大部分专家在特定任务下几乎不会被激活，因此 CPU offload 的实际影响远小于理论上的"跨 PCIe 传输全部权重"。

\textbf{实测潜力}：KTransformers 宣称可以在 4$\times$RTX 4090（24 GB $\times$ 4）的配置上运行 DeepSeek-V3（671B MoE），实现 10--15 tokens/s 的可用速度。对于 GLM-5.1 的类似架构，理论上也可以在 4$\times$A6000（48 GB $\times$ 4 = 192 GB）上实现可用部署。但这一方案尚处实验阶段，生产稳定性有待验证。

\subsection{xLLM（华为昇腾方案）}

对于有信创/国产化需求的团队，xLLM 是基于华为昇腾（Ascend）硬件的推理框架。
\begin{itemize}
    \item 需 Ascend 910B 或更新型号
    \item 生态远不如 CUDA 成熟，模型适配需额外工作量
    \item 适合信创合规要求下的国企/政府场景
\end{itemize}

% ======================================================================
% 第五章：成本预估
% ======================================================================

\section{成本预估}

\subsection{硬件采购成本}

\begin{table}[htbp]
\centering
\caption{GLM-5.1 各部署方案的硬件采购成本估算（2026 年 5 月参考价，已修正）}
\begin{tabular}{p{3.5cm}p{1.8cm}p{2.3cm}p{3cm}p{2cm}}
\hline
\textbf{方案} & \textbf{GPU 配置} & \textbf{GPU 总价} & \textbf{服务器总价} & \textbf{总预算} \\
\hline
A: BF16 全精度 & 16$\times$H200 SXM & \$560K & \$150--200K & \$710--760K \\
A: BF16 全精度 & 24$\times$H100 SXM & \$672K & \$180--250K & \$852--922K \\
B-1: FP8 官方 & 8$\times$H200 SXM & \$280K & \$80--120K & \$360--400K \\
B-2: FP8 H100扩容 & 16$\times$H100 SXM & \$448K & \$120--180K & \$568--628K \\
B-3: FP8 H100最低 & 12$\times$H100 SXM & \$336K & \$100--150K & \$436--486K \\
C: Q4 GGUF & 8$\times$A100 SXM & \$120K & \$50--80K & \$170--200K \\
C: Q4 GGUF & 12$\times$A6000+ & \$82K & \$60--90K & \$142--172K \\
D: NVFP4 & 8$\times$B200 & \$280K & \$100--150K & \$380--430K \\
D: MXFP4 & 8$\times$MI300X & \$120K & \$80--120K & \$200--240K \\
\hline
\end{tabular}
\end{table}

\textbf{说明}：服务器总价包括双路服务器平台（CPU、内存 512 GB+、NVMe 存储、InfiniBand 网络、散热、电源等）。多节点方案（16/24 卡）需要 2--3 台 8-GPU 服务器通过 InfiniBand 互联，服务器成本显著增加。GPU 价格为 2026 年 5 月大致的市场参考价。

\subsection{云服务成本对比}

\subsubsection{按需实例 TCO 分析}

以 8$\times$H200 的 HuggingFace Inference Endpoint 为例（官方 recipe 配置）：

\begin{itemize}
    \item \textbf{时租成本}：约 \$40--55/小时（8$\times$H200 GPU 实例 + HuggingFace 平台溢价）
    \item \textbf{月租成本}（720 小时/月）：\$45 $\times$ 720 = \$32,400/月
    \item \textbf{年租成本}：约 \$388,800/年（无预留折扣）
    \item \textbf{3 年 TCO}：约 \$1,166,400
\end{itemize}

以 8$\times$H100 的云实例（仅能跑 Q4 GGUF，不可跑 FP8）为例：

\begin{itemize}
    \item \textbf{时租成本}：约 \$25--30/小时
    \item \textbf{月租成本}：\$28 $\times$ 720 = \$20,160/月
    \item \textbf{3 年 TCO}：约 \$725,760
\end{itemize}

\subsubsection{API 调用成本}

以 Together AI API 的典型 GLM-5.1 定价为例（参考同级别 MoE 模型）：

\begin{itemize}
    \item \textbf{Input tokens}：约 \$3--6 / 百万 tokens
    \item \textbf{Output tokens}：约 \$8--15 / 百万 tokens
    \item \textbf{按典型 Agent 场景估算}：单次 Agent 任务平均消耗 50K input + 10K output tokens，成本约 \$0.15--0.45/次
    \item \textbf{如果每天 100 次 Agent 调用}：\$15--45/天，\$450--1,350/月
    \item \textbf{如果每天 1,000 次 Agent 调用}：\$150--450/天，\$4,500--13,500/月
\end{itemize}

\subsubsection{自有硬件 vs 云服务的盈亏平衡点}

\begin{table}[htbp]
\centering
\caption{自有硬件 vs 云服务/API 的 3 年 TCO 对比（单位：万美元，已修正）}
\begin{tabular}{p{3.5cm}p{2cm}p{2cm}p{2cm}p{2.5cm}}
\hline
\textbf{方案} & \textbf{初始投入} & \textbf{年运维} & \textbf{3 年 TCO} & \textbf{适用场景} \\
\hline
自购 8$\times$H200 (FP8) & \$38 & \$4 & \$50 & 7$\times$24 满负荷使用 \\
自购 8$\times$A100 (Q4) & \$18 & \$2.5 & \$25.5 & 中等负载+低精度 \\
云 GPU 8$\times$H200 时租 & \$0 & \$20--39/年 & \$60--117 & 间歇性使用（50\%）\\
API 按量计费 & \$0 & \$5--16/年 & \$15--48 & 低频 Agent 场景 \\
\hline
\end{tabular}
\end{table}

\textbf{粗略决策规则}：
\begin{itemize}
    \item 每天使用少于 8 小时 $\rightarrow$ 云 GPU 按时租更划算
    \item 每天使用超过 16 小时、负载可预测 $\rightarrow$ 自购硬件更划算
    \item 低频 Agent 调用（每天 $<$ 500 次）$\rightarrow$ API 按量计费最经济
    \item 高频调用 + 对延迟敏感 + 涉及敏感数据 $\rightarrow$ 必须自购硬件
\end{itemize}

\subsection{电力与运维成本}

\begin{itemize}
    \item 8$\times$H200 服务器功耗：约 7,000--9,000W（含 CPU、内存、散热、InfiniBand 交换）
    \item 年电力成本（按 \$0.12/kWh）：8 kW $\times$ 8,760 h $\times\$0.12 \approx$ \$8,410/年
    \item 16$\times$H100 方案（2 节点）：功耗约 11,000--14,000W，年电力成本约 \$11,500--14,700
    \item 数据中心托管（8U 机柜 + 冗余电力 + InfiniBand 网络）：约 \$2,000--3,500/月
    \item 运维人力：通常由现有工程团队兼任，边际成本低
\end{itemize}

% ======================================================================
% 第六章：性能预期与实测分析
% ======================================================================

\section{性能预期与实测分析}

\subsection{各方案理论吞吐量估算}

Decode 阶段的理论吞吐量可以通过"显存带宽 $\div$ 模型读取量"公式估算。对于 MoE 模型，每次生成一个 token 不需要读取全部 744B 权重（仅激活 40B），但 KV Cache 的读取仍依赖于上下文长度。

\begin{table}[htbp]
\centering
\caption{各方案的预期 Decode 速度估算（单请求，32K 上下文）}
\begin{tabular}{p{3cm}p{1.8cm}p{2.2cm}p{2.5cm}p{2cm}}
\hline
\textbf{方案} & \textbf{GPU} & \textbf{每 token 数据读取量} & \textbf{有效带宽} & \textbf{估算速度} \\
\hline
A: BF16 & 16$\times$H200 & $\sim$90 GB & 16$\times$4.8 TB/s & 50--80 tok/s \\
B-1: FP8 & 8$\times$H200 & $\sim$48 GB & 8$\times$4.8 TB/s & 45--65 tok/s \\
B-2: FP8 & 16$\times$H100 & $\sim$25 GB & 16$\times$3.35 TB/s & 40--55 tok/s \\
B-3: FP8 & 12$\times$H100 & $\sim$33 GB & 12$\times$3.35 TB/s & 35--50 tok/s \\
C: Q4 GGUF & 8$\times$A100 & $\sim$35 GB & 8$\times$2.0 TB/s & 20--35 tok/s \\
C: Q4 GGUF & 12$\times$A6000 + CPU & $\sim$25 GB+CPU & 12$\times$0.96 + CPU & 5--10 tok/s \\
\hline
\end{tabular}
\end{table}

\subsection{上下文长度对显存的影响}

GLM-5.1 支持 128K 上下文，但实际可用上下文受 GPU 显存约束。

\begin{table}[htbp]
\centering
\caption{不同上下文长度下 KV Cache 的额外显存占用（FP8, 8$\times$H200 官方配置）}
\begin{tabular}{p{2.5cm}p{2.2cm}p{3cm}p{2.2cm}p{2.2cm}}
\hline
\textbf{上下文} & \textbf{KV Cache/GPU} & \textbf{权重+KV+开销} & \textbf{141GB 占用率} & \textbf{状态} \\
\hline
4K & $\sim$0.3 GB & $\sim$97 GB & 69\% & 宽松 \\
8K & $\sim$0.5 GB & $\sim$97 GB & 69\% & 宽松 \\
32K & $\sim$2.0 GB & $\sim$99 GB & 70\% & 正常 \\
128K & $\sim$6.0 GB & $\sim$103 GB & 73\% & 正常 \\
200K & $\sim$9.0 GB & $\sim$106 GB & 75\% & 正常 \\
\hline
\end{tabular}
\end{table}

\textbf{说明}：得益于 MLA（KV LoRA rank=512，压缩比约 93\%）和 DSA 稀疏注意力（每查询仅取 top-2048 个 KV），GLM-5.1 的 KV Cache 远小于传统 MHA/GQA 架构。在 8$\times$H200 官方配置下，即使 200K 上下文单请求的 KV Cache 也仅约 9 GB/卡。\textbf{这与传统模型动辄数十 GB 的 KV Cache 形成鲜明对比，是 GLM-5.1 架构优势的集中体现。}

\textbf{缓解措施}：
\begin{itemize}
    \item vLLM 的 FP8 KV Cache 量化：将 KV Cache 从 FP16 压缩到 FP8，减少约 50\% 的 KV Cache 显存
    \item Prefix Caching：在 Agent 场景中，多轮对话共享 system prompt 和历史前缀，避免重复存储
    \item vLLM 的 Automatic Prefix Caching（APC）：自动识别和复用公共前缀
\end{itemize}

\subsection{并发能力评估}

\begin{table}[htbp]
\centering
\caption{不同配置下的预估并发能力（32K 上下文）}
\begin{tabular}{p{3cm}p{2.5cm}p{2.5cm}p{2.5cm}p{2.2cm}}
\hline
\textbf{配置} & \textbf{可用总显存} & \textbf{单请求 KV Cache} & \textbf{理论并发} & \textbf{实际并发} \\
\hline
8$\times$H200 (FP8) & $\sim$1,000 GB & $\sim$2 GB/req & $\sim$500 & 10--20 \\
16$\times$H100 (FP8) & $\sim$1,100 GB & $\sim$2 GB/req & $\sim$550 & 8--15 \\
12$\times$H100 (FP8) & $\sim$850 GB & $\sim$2 GB/req & $\sim$425 & 5--8 \\
8$\times$A100 (Q4 GGUF) & $\sim$550 GB & $\sim$2 GB/req & $\sim$275 & 3--5 \\
\hline
\end{tabular}
\end{table}

\textbf{说明}：理论并发受 KV Cache 容量约束，实际并发还受 GPU 算力和 Decode 带宽约束。对于 Agent 场景（长会话、频繁工具调用），建议保守估计并发数。

% ======================================================================
% 第七章：推荐方案矩阵
% ======================================================================

\section{推荐方案矩阵}

\subsection{按场景推荐}

\begin{table}[htbp]
\centering
\caption{GLM-5.1 部署方案场景推荐矩阵（修正版）}
\begin{tabular}{p{2.5cm}p{3.2cm}p{2.8cm}p{3.5cm}}
\hline
\textbf{场景} & \textbf{推荐方案} & \textbf{预算范围} & \textbf{备注} \\
\hline
个人研究者 & API + HuggingChat & \$50--500/月 & 无需硬件，HuggingChat 免费体验 \\
小型创业团队 & 云 GPU 8$\times$H200 时租 + API 混合 & \$2K--15K/月 & 弹性伸缩，按需付费 \\
中型研发团队 & 自购 8$\times$A100 + Q4 GGUF & \$18--25 万 & 最低自建方案，一次性投入 \\
大型研发团队 & 自购 8$\times$H200 + FP8 & \$36--40 万 & 官方 recipe 标准配置 \\
企业级平台 & 自购 16$\times$H200 + BF16 & \$71--76 万 & 全精度、高并发、低延迟 \\
信创合规 & xLLM + 华为昇腾 & 视配置 & 国产化替代，生态受限 \\
\hline
\end{tabular}
\end{table}

\subsection{决策流程图}

以下是一个简化的决策逻辑：

\begin{enumerate}
    \item \textbf{是否有 \$36 万以上的硬件预算？}
    \begin{itemize}
        \item 是 $\rightarrow$ 方案 B-1（自购 8$\times$H200 + FP8），官方标准配置
        \item 否 $\rightarrow$ 进入问题 2
    \end{itemize}
    \item \textbf{是否有 \$18 万以上的硬件预算？}
    \begin{itemize}
        \item 是 $\rightarrow$ 方案 C（自购 8$\times$A100 + Q4 GGUF），最低自建方案
        \item 否 $\rightarrow$ 进入问题 3
    \end{itemize}
    \item \textbf{日调用量是否超过 500 次 Agent 任务？}
    \begin{itemize}
        \item 是 $\rightarrow$ 云 GPU 按时租（8$\times$H200 约 \$40--55/时），弹性调配
        \item 否 $\rightarrow$ 进入问题 4
    \end{itemize}
    \item \textbf{是否有数据安全/合规要求（不可出内网）？}
    \begin{itemize}
        \item 是 $\rightarrow$ KTransformers + 4$\times$A6000 + CPU offload（实验性最低自建方案）
        \item 否 $\rightarrow$ API 按量计费（Together AI / HuggingFace），零运维
    \end{itemize}
\end{enumerate}

\subsection{最终建议}

\begin{personalnote}
\textbf{对于大多数研发团队}，GLM-5.1 的部署建议如下（已根据显存物理上限修正）：

\begin{enumerate}
    \item \textbf{PoC/探索阶段}：先使用 HuggingChat（免费）或 Together AI API（$<\$500$/月）体验模型能力，确认是否满足业务需求。\textbf{不要急于采购硬件}。

    \item \textbf{小规模使用阶段}：使用云 GPU 按时租（RunPod/Vast.ai 8$\times$A100 $\approx\$12--20$/时用于 Q4 GGUF；8$\times$H200 $\approx\$40--55$/时用于 FP8 官方 recipe）。每周运行 20--40 小时，月成本约 \$1,000--9,000。此时无需任何硬件投入。

    \item \textbf{规模化阶段}（确认 GLM-5.1 为团队核心生产力工具）：
    \begin{itemize}
        \item \textbf{推荐方案}：采购 8$\times$A100 SXM 服务器（约 \$18--20 万），部署 Q4\_K\_M GGUF 量化模型。3 年 TCO 约 \$25 万，对比 8$\times$H200 云租方案可节省 50\%+。
        \item \textbf{高性能方案}：采购 8$\times$H200 SXM 服务器（约 \$36--40 万），部署官方 FP8 模型。这是当前 GLM-5.1 的最佳部署配置。
        \item \textbf{注意}：\textbf{不要购买 8$\times$H100 用于 FP8 部署}——93 GB/卡的权重超过 H100 80 GB 物理上限。H100 仅适用于 Q4 GGUF 或更大规模的集群。
    \end{itemize}

    \item \textbf{最终结论}：当前 \textbf{GLM-5.1 的硬件门槛远超消费级/个人开发者预算}。最低自建方案为 8$\times$A100 + Q4 GGUF（约 \$18 万），FP8 官方 recipe 需 8$\times$H200（约 \$36 万）。对于个人开发者和小团队，建议直接使用云 API，或转向更轻量但同样优秀的模型（如 Qwen3-Coder-32B、DeepSeek-Coder-V2，均可在单卡或双卡上高效运行）。
\end{enumerate}
\end{personalnote}

% ======================================================================
% 附录：快速参考
% ======================================================================

\section*{附录 A：关键数据速查}

\begin{table}[htbp]
\centering
\caption{GLM-5.1 关键指标速查表（已修正）}
\begin{tabular}{p{3.8cm}p{8cm}}
\hline
\textbf{指标} & \textbf{数值} \\
\hline
总参数量 & 744B（78 层，75 层 MoE）\\
激活参数量 & 36B (4.8\%) \\
上下文窗口 & 200K tokens（max\_position\_embeddings=202752）\\
原始精度 & BF16 ($\sim$1,488 GB) \\
官方推荐精度 & FP8 ($\sim$744 GB) \\
最低可用量化 & Q4\_K\_M ($\sim$465 GB) \\
官方推荐 GPU & \textbf{8$\times$ H200 SXM (141 GB)} ——注意不是 H100！\\
FP8 最低 GPU 配置（H100）& \textbf{12$\times$ H100 SXM (80 GB)}，临界；推荐 16 张 \\
Q4 GGUF 最低 GPU 配置 & 8$\times$ A100 SXM (80 GB) \\
实验性最低方案 & KTransformers + 4$\times$ A6000 (48 GB) + 512 GB RAM \\
推理框架（推荐） & vLLM v0.20.2\textsuperscript{[1]} / SGLang v0.5.11\textsuperscript{[1]} \\
推理框架（量化） & llama.cpp / KTransformers \\
API 服务商 & Together AI, HuggingFace Inference \\
免费体验 & HuggingChat (chat.z.ai) \\
开源协议 & MIT / Apache 2.0 \\
\hline
\end{tabular}
\end{table}

\section*{附录 B：社区量化资源链接}

\begin{itemize}
    \item 官方模型：\href{https://huggingface.co/zai-org/GLM-5.1}{zai-org/GLM-5.1}（BF16）| \href{https://huggingface.co/zai-org/GLM-5.1-FP8}{zai-org/GLM-5.1-FP8}
    \item GGUF 量化（推荐）：\href{https://huggingface.co/unsloth/GLM-5.1-GGUF}{unsloth/GLM-5.1-GGUF} | \href{https://huggingface.co/bartowski/zai-org_GLM-5.1-GGUF}{bartowski/GLM-5.1-GGUF}
    \item NVFP4：\href{https://huggingface.co/lukealonso/GLM-5.1-NVFP4}{lukealonso/GLM-5.1-NVFP4}
    \item MXFP4：\href{https://huggingface.co/amd/GLM-5.1-MXFP4}{amd/GLM-5.1-MXFP4}
    \item MLX（Apple Silicon）：\href{https://huggingface.co/mlx-community/GLM-5.1-DQ4plus-q8}{mlx-community/GLM-5.1-DQ4plus-q8}
    \item AWQ 4-bit：\href{https://huggingface.co/cyankiwi/GLM-5.1-AWQ-4bit}{cyankiwi/GLM-5.1-AWQ-4bit}
    \item REAP 剪枝+量化：\href{https://huggingface.co/0xSero/GLM-5.1-478B-A42B-REAP-NVFP4}{0xSero/GLM-5.1-REAP-NVFP4}
    \item 官方 GitHub：\href{https://github.com/zai-org/GLM-5}{github.com/zai-org/GLM-5}
    \item 技术报告：\href{https://arxiv.org/abs/2602.15763}{arXiv:2602.15763}
\end{itemize}

% ======================================================================
% 参考文献
% ======================================================================

\section*{参考文献与数据来源}

\begin{enumerate}
    \item \textbf{GLM-5 官方 GitHub 仓库}（Apache 2.0）。包含 GLM-5 与 GLM-5.1 的部署命令、框架版本要求、模型参数说明。\newline
    \url{https://github.com/zai-org/GLM-5}（2026-05-20 访问）

    \item \textbf{GLM-5.1 HuggingFace 模型页}（MIT）。包含模型信息、推理代码示例、量化变体链接、社区讨论。\newline
    \url{https://huggingface.co/zai-org/GLM-5.1}（2026-05-20 访问）

    \item \textbf{vLLM GLM-5.1 Recipe}（vLLM 官方部署指南）。\textbf{明确标注目标硬件为 "8xH200 (or H20) GPUs (141GB$\times$8)"}，包含完整 vLLM 部署命令和性能基准测试方法。这是本文硬件需求分析的关键来源。\newline
    \url{https://github.com/vllm-project/recipes/blob/main/GLM/GLM5.md}（2026-05-20 访问）

    \item \textbf{GLM-5.1 模型配置文件}（config.json）。包含全部架构超参数：78 层、256 路由专家、MLA/KV LoRA/Indexer 配置、200K 上下文窗口。\newline
    \url{https://huggingface.co/zai-org/GLM-5.1/blob/main/config.json}（2026-05-20 访问）

    \item \textbf{GLM-5 技术报告}。Z.ai 团队。\textit{GLM-5: from Vibe Coding to Agentic Engineering}。arXiv:2602.15763。\newline
    \url{https://arxiv.org/abs/2602.15763}

    \item \textbf{GLM-5.1 官方 FP8 量化模型}。\newline
    \url{https://huggingface.co/zai-org/GLM-5.1-FP8}

    \item \textbf{Unsloth GLM-5.1 GGUF 量化}（社区）。包含完整量化变体（UD-Q2\_K 至 UD-Q8\_0、IQ 系列），总仓库大小 11.1 TB。\newline
    \url{https://huggingface.co/unsloth/GLM-5.1-GGUF}

    \item \textbf{Bartowski GLM-5.1 GGUF 量化}（社区，备选）。\newline
    \url{https://huggingface.co/bartowski/zai-org_GLM-5.1-GGUF}

    \item \textbf{SGLang GLM-5.1 Cookbook}。\newline
    \url{https://cookbook.sglang.io/autoregressive/GLM/GLM-5.1}

    \item \textbf{KTransformers GLM-5.1 教程}。\newline
    \url{https://github.com/kvcache-ai/ktransformers/blob/main/doc/en/kt-kernel/GLM-5.1-Tutorial.md}
\end{enumerate}

\textbf{数据一致性说明}：
\begin{itemize}
    \item 总参数量：GitHub 官方 README 表述为 744B\textsuperscript{[1]}；HuggingFace 页面自动统计显示为 754B\textsuperscript{[2]}。差异约 1.3\%，可能源于参数计数方法（是否包含 embedding/LM head 等固定参数的不同处理）。本文以 GitHub 官方数字 744B 为准进行显存计算。
    \item 激活参数量：GitHub 官方 README 表述为 40B\textsuperscript{[1]}。从 config.json 架构参数推算约为 40--49B（取决于是否计入 Attention/Indexer/Embedding 等非 MoE 部分的完整激活）。
    \item 硬件要求：GitHub README 的部署命令\textbf{未指定 GPU 型号}\textsuperscript{[1]}——仅给出 \texttt{--tensor-parallel-size 8} 和 \texttt{--gpu-memory-utilization 0.85}。vLLM recipe 页面\textbf{明确标注} "Serving FP8 Model on 8xH200 (or H20) GPUs (141GB$\times$8)"\textsuperscript{[3]}。HuggingFace 页面命令完全省略了 TP 配置\textsuperscript{[2]}。\textbf{三个来源的详细程度差异是初版产生"8$\times$H100 可行"这一错误判断的直接原因。}
\end{itemize}

\end{document}
